% TOP_PROBLEM: {"id":30007533,"problem_number":"LOCAL-30007533","title":"Public-investment returns","category_id":21,"category":{"id":21,"name":"economics","display_name":"Economics","description":"Open economic theory statements, published research agendas, and proposed scoped specifications; question provenance and historical priority are qualified per record.","slug":"economics","order_index":21,"created_at":"2026-10-08T00:00:00Z"},"status":"provenance_pending","external_url":null,"legacy_ids":[],"published":false,"statement":"Identify the distribution of public-project net social returns including productivity effects, delays, maintenance, spillovers, and financing costs.","economics":{"original_id":22,"field":"Fiscal policy and sovereign debt","kind":"identification","formalization_basis":"proposed_specification","source_status":"not_verified","priority_status":"not_established","priority_note":"No explicit primary open-problem statement matching this catalogue item was verified; supporting research is not evidence of first formulation.","earliest_verified_reference":null,"current_reference":{"authors":["International Monetary Fund"],"title":"Fiscal Monitor, October 2025: Spending Smarter—How Efficient and Well-Allocated Public Spending Can Boost Economic Growth","year":2025,"url":"https://www.elibrary.imf.org/abstract/book/9798229024938/CH001.xml","doi":null,"locator":"Chapter 1, primary publisher summary/abstract","evidence_summary":"The report assesses expenditure allocation, efficiency, and growth benefits. It offers evidence and recommendations; a precise unresolved project-return estimand matching the original catalogue item was not verified."},"formalization_note":"New project-level identification target. The current IMF report motivates efficiency and allocation but no matching explicit open statement was verified."},"statusQualification":"Provenance pending: an explicit published statement of this exact open question was not verified. This is a catalogue-authored candidate specification, not a certified open conjecture. New project-level identification target. The current IMF report motivates efficiency and allocation but no matching explicit open statement was verified.","statusReviewedAt":"2026-10-08"}
% FIRST_POSTED: 2026-10-08
% ENTRY_KIND: statement_only
% !TeX program = lualatex
\documentclass[11pt]{article}
\usepackage{fontspec}
\usepackage[a4paper,margin=25mm]{geometry}
\usepackage{amsmath,amssymb,mathtools}
\usepackage{xurl}
\usepackage[hidelinks]{hyperref}
\newcommand{\sref}[2]{\hyperlink{source:#1:#2}{[S#2]}}
\setlength{\emergencystretch}{3em}
\begin{document}
% problemId:30007533
\section*{Public-investment returns}\label{30007533}
\subsection{Definitions and mathematical statement}
Fix a class of public infrastructure projects \(j\) with construction cost \(K_j\), maintenance stream \(M_{jt}\), completion delay \(d_j\), and independently measured service capacity. Let \(A_{jt}(1)\) and \(A_{jt}(0)\) denote private-sector productivity with and without project implementation. Let \(S_{jt}(A)\) be the gross social value of output and access at productivity level \(A\), net of congestion and environmental damages; its valuation assumptions must be declared. At discount rate \(r>0\), define
\[
V_j=E\sum_{t\ge d_j}(1+r)^{-t}
[S_{jt}(A_{jt}(1))-S_{jt}(A_{jt}(0))-M_{jt}]
-K_j-C_j^F,
\]
where \(C_j^F\) is the welfare cost of the specified financing scheme, including displaced expenditures or distortionary taxes.
Identify the distribution of \(V_j\) and productivity gains using credible project assignment, timing, or funding instruments. Address endogenous project location, anticipation, spillovers to untreated areas, and selective completion. Separate physical infrastructure productivity from utilization and administrative efficiency, and quantify the consequences of delays and cost overruns. Establish which measured project characteristics predict positive net value out of sample. A country-level spending regression cannot by itself identify project returns or financing costs. This is a scoped project-evaluation problem; the general possibility that productive public investment benefits growth is already supported by substantial research.

\par\textbf{Formulation and scope.} New project-level identification target. The current IMF report motivates efficiency and allocation but no matching explicit open statement was verified.

\par\textbf{Open-question provenance.} An explicit published open-question statement was not verified. Sources below are supporting literature and must not be treated as a first listing.

\par\textbf{Historical priority.} No explicit primary open-problem statement matching this catalogue item was verified; supporting research is not evidence of first formulation.

\subsection{Short English statement}
Identify the distribution of public-project net social returns including productivity effects, delays, maintenance, spillovers, and financing costs.

\subsection{Sources}
\begin{itemize}\raggedright
\item[\textnormal{[S1]}]\hypertarget{source:30007533:1}{} International Monetary Fund. 2025. Fiscal Monitor, October 2025: Spending Smarter—How Efficient and Well-Allocated Public Spending Can Boost Economic Growth. Chapter 1, primary publisher summary/abstract.
\url{https://www.elibrary.imf.org/abstract/book/9798229024938/CH001.xml}.
{\small\textit{Supporting or current literature; see evidence qualification. The report assesses expenditure allocation, efficiency, and growth benefits. It offers evidence and recommendations; a precise unresolved project-return estimand matching the original catalogue item was not verified.}}
\end{itemize}
\end{document}
